\section{Introduction}\label{sec:introduction}

Grid operators need reliable day-ahead views of where load will rise and how tightly the system is constrained. In developing economies, rapid demand growth often coincides with supply limits and load shedding, so planners cannot rely on demand forecasts alone; they also need a readable signal of operational pressure. Bangladesh's national power system reports daily data for nine administrative divisions, with evening-peak demand appearing alongside limitation, reserve, and shedding fields in the same records \cite{Hossain2025BangladeshSmartGrid}. Planning is therefore inherently spatial and system-wide at once.

In practice, however, load forecasting and stress monitoring are still commonly built as separate tools. Classical and ensemble models \cite{Breiman2001RandomForest,Chen2016XGBoost,Musbah2025LoadForecasting} predict demand without modelling links between divisions. Shallow recurrent or graph-convolution networks \cite{Hochreiter1997LSTM,Zhao2020TGCN} may struggle to combine a full weekly lookback with informative cross-node structure. Graph methods developed for residential or microgrid settings often assume geographical neighbours \cite{Kipf2017GCN}, which may not match how regional loads actually move together on a national footprint. Operational stress---driven by reserve shortfall, limitation pressure, and shedding---is rarely forecast jointly with regional demand; most shedding work \cite{Paphitis2025SpatialLoadShedding} targets control rather than next-day stress indices. Explainability is also uncommon in graph-temporal models that output both load and stress \cite{Baur2024XAIReview,Lundberg2017SHAP}.

We treat day-ahead forecasting as a supervised multi-task problem \cite{Caruana1997Multitask} over nine divisions. Task~1 predicts next-day regional evening-peak demand in MW. Task~2 predicts a graph-level Operational Stress Index (OSI) in $[0,1]$ that summarises shedding intensity, reserve margin, and limitation pressure. Demand and OSI are related in operations but not redundant: demand describes expected load by region, while OSI captures constraint pressure that demand alone does not fully explain. The question we ask is whether one spatio-temporal model can learn both outputs from historical Bangladesh grid records under strict chronological evaluation.

Our approach is the Correlation-Aware Multi-Task Forecasting Framework, implemented as a Parallel-Fusion Spatio-Temporal Graph Transformer (PF-STGT). The model takes a seven-day lookback of division-level and national features together with a correlation graph estimated from training data, combines graph-aware spatial attention with temporal transformers \cite{Vaswani2017Attention}, and decodes regional demand and bounded OSI through separate heads. Training uses a repaired multi-task loss so neither task is drowned out during optimisation. We evaluate against classical, recurrent, and graph-convolution baselines, run configuration ablations, apply Wilcoxon testing with Bonferroni correction, and report post-hoc attributions for the selected model.

This paper makes five contributions: (i)~a correlation-graph PF-STGT for joint day-ahead demand and OSI forecasting on Bangladesh national data; (ii)~evidence that correlation-only adjacency outperforms geographical and hybrid priors on this dataset; (iii)~a practical multi-task loss calibration that stabilises joint training; (iv)~a full benchmark and ablation programme with inferential testing; and (v)~integrated explainability with multi-method attribution for both tasks.

Section~\ref{sec:methodology} describes the method; Sections~\ref{sec:experimental_setup}--\ref{sec:results} present experiments and results; Sections~\ref{sec:discussion}--\ref{sec:conclusion} interpret the findings and conclude.
